\documentclass[12pt, letterpaper]{article}
\usepackage[margin=1.2in]{geometry}
\usepackage{ebgaramond}
\usepackage[T1]{fontenc}
\usepackage{microtype}
\usepackage{setspace}
\usepackage{parskip}
\usepackage{titling}
\usepackage{fancyhdr}

\setlength{\parindent}{2em}
\onehalfspacing

\pagestyle{fancy}
\fancyhf{}
\fancyhead[L]{\small\textit{The Artificial Aurora}}
\fancyhead[R]{\small\textit{NASA Mission Series v1.10}}
\fancyfoot[C]{\small\thepage}
\renewcommand{\headrulewidth}{0.4pt}

\title{\Huge\textbf{The Artificial Aurora}\\[0.5em]\large\textit{A short story from Mission 1.10 --- Explorer 4}}
\author{NASA Mission Series\\[0.3em]\small Miles Tiberius Foxglove --- Mukilteo, WA}
\date{March 30, 2026}

\begin{document}
\maketitle
\thispagestyle{empty}

\vspace{1em}
\noindent\rule{\textwidth}{0.4pt}
\vspace{1em}

\noindent The South Atlantic was flat and black and without feature. Seaman Apprentice David Kohler stood on the weather deck of USS Norton Sound and looked at water that had no business being so calm. The ship had been at general quarters for twenty minutes, which meant the hatches were dogged, the ventilation was off, and eighty-seven men were sweating in compartments sealed like cans. Kohler was one of the lucky ones --- a topside watch, a pair of welding goggles pushed up on his forehead, and a clipboard with a form titled OPTICAL OBSERVATION REPORT that he had been told to fill out when the thing happened. Nobody had told him what the thing was.

It was August 27, 1958, and the ship was a thousand miles from anywhere. They had sailed from San Diego six weeks ago, down the coast of South America, around into the Atlantic, and stopped here --- 38 degrees south, 11 degrees west --- in water so deep that the anchor would have fallen for five miles before it touched mud. The captain had said they were conducting a scientific experiment. The executive officer had said they were testing a weapon. The chief said shut up and do your job. Kohler was nineteen years old and did what the chief said.

The X-17A sat on its launcher amidships, a three-stage rocket fifty-six feet tall, painted white, aimed at the sky. Kohler had watched them load it three days ago --- a long cylinder with a warhead smaller than a garbage can. One point seven kilotons. The physics team from the Defense Nuclear Agency called it a ``device.'' The weapon was already aboard; the word was just a courtesy.

\begin{center}
$\bullet\quad\bullet\quad\bullet$
\end{center}

At 0228 UTC the launcher belched white smoke and the X-17A left the ship with a sound like a building falling. Kohler pulled the welding goggles down over his eyes and watched the rocket climb --- a white flame shrinking to a star, the star shrinking to a point, the point climbing east and upward into the absolute dark of the South Atlantic night sky. Around him, the ship trembled with the backwash of the first stage, then went quiet. The rocket was beyond hearing now. It was climbing through the mesosphere, through the thermosphere, up to two hundred kilometers where the atmosphere gives way to the void. Up to where the Van Allen belts begin.

He counted. He had been told to count. At one hundred and twelve seconds, the sky turned white.

Not the entire sky. A disk of light, directly overhead, maybe ten degrees across, appeared where the point of flame had been. It was brighter than anything Kohler had ever seen through welding glass --- brighter than the arc welding he had done in A-school, brighter than the flashbulbs the photographers used at ship's events. It was the color of the inside of a star: not yellow, not orange, not blue, but all colors compressed into a single white that contained everything and disclosed nothing. It lasted less than a second at full intensity, then expanded and dimmed, spreading outward like a ripple in a pond made of light.

And then the aurora.

Kohler had never seen an aurora. He was from Modesto, California, where the night sky was orange from the streetlights and the most dramatic thing overhead was the contrail from Travis Air Force Base. But he had read about aurora borealis in a National Geographic his mother kept in the living room, and what he saw now matched the photographs --- except it was in the wrong ocean, at the wrong latitude, and it was caused by something he had helped load onto a launcher three days ago.

Green. The light was green. Not the green of grass or of traffic lights or of the starboard running light, but a green that lived in a register Kohler had no name for --- electric, vivid, slightly blue at the edges, shifting and pulsing like a curtain hung from nothing. It spread in a band across the sky from southeast to northwest, following a curve that corresponded to nothing on the horizon. The band was maybe thirty degrees wide and it flickered with an internal rhythm, as if something behind it was breathing.

\begin{center}
$\bullet\quad\bullet\quad\bullet$
\end{center}

What Kohler was seeing --- what eighty-seven men on the Norton Sound and the crews of nine other ships and two aircraft were simultaneously recording on film, spectroscopes, magnetometers, and clipboards identical to his --- was the Christofilos effect. A Greek physicist named Nicholas Christofilos, working at the Lawrence Radiation Laboratory in Livermore, had predicted in 1957 that a nuclear detonation at high altitude would inject a shell of fission electrons into Earth's magnetic field. The electrons, being charged, would spiral along the magnetic field lines, trapped between mirror points near the poles. They would form an artificial radiation belt --- a thin, hot, invisible shell of trapped particles wrapping the planet at the altitude of the explosion. And where the trapped electrons descended toward the mirror points, some would collide with atmospheric molecules and produce light. Auroral light. Artificial aurora.

Christofilos had been right about everything. The electrons spiraled. The belt formed. The aurora glowed. Eight hundred kilometers above the Norton Sound, a satellite called Explorer 4 --- launched four weeks earlier specifically to carry scintillation counters through this experiment --- was recording the arrival of the trapped electrons. Its counters, which had been ticking quietly at the background rate of natural cosmic radiation, surged. The artificial belt was real. It was measurable. It was exactly where Christofilos said it would be.

Kohler did not know any of this. He was filling out his clipboard. Under INTENSITY he wrote ``very bright.'' Under DURATION he wrote ``approx 5 min visible.'' Under COLOR he wrote ``green, some blue at edges.'' Under ADDITIONAL OBSERVATIONS he hesitated, then wrote: ``Like a curtain. Moving. Not natural.''

\begin{center}
$\bullet\quad\bullet\quad\bullet$
\end{center}

They did it two more times. Argus II on August 30, at 240 kilometers. Argus III on September 6, at 540 kilometers --- the highest, the deepest injection into the trapping region. Each time the flash, the expanding shell, the aurora. Each time Kohler stood on the weather deck with his goggles and his clipboard and wrote down what he saw. Each time Explorer 4 recorded the surge of particles. Each time the artificial belt glowed for a while and then faded, the electrons scattering into the loss cone and raining down into the atmosphere where they produced one last flicker of light before being absorbed.

The third test was the most beautiful. At 540 kilometers, the injection was deep enough that the belt persisted for weeks rather than days. The aurora from Argus III was visible from Apia, Samoa, where a tracking station reported green bands across the sky that puzzled local observers who had never seen aurora at fourteen degrees south latitude. In the Azores, the Royal Observatory photographed green streaks that they initially attributed to instrumentation error. In Argentina, a farmer reported green lights to a newspaper that printed the story next to a report on flying saucers. Nobody told any of them what it was. The project was classified. Twenty-six hundred people knew. The rest of the world saw green light in the wrong part of the sky and drew their own conclusions.

\begin{center}
$\bullet\quad\bullet\quad\bullet$
\end{center}

Three thousand miles north of the Norton Sound and twelve thousand feet above sea level, the tundra was patient. On the high plateau of the Brooks Range in Alaska, where the permafrost goes down a thousand feet and the growing season lasts six weeks, a lichen called \textit{Cladonia stellaris} was doing what it always does: growing at two millimeters per year, absorbing moisture from the air, filtering whatever the sky chose to deliver.

\textit{Cladonia stellaris} is a pale, yellow-green organism that looks like a tiny forest of branching coral. Each branch splits into finer branches, and each fine branch ends in a star-shaped tip --- \textit{stellaris}, star-tipped. A single colony can be fifteen centimeters across and fifty years old. The lichen has no roots. It has no vascular system. It has no way to reject what falls on it. It absorbs indiscriminately: water, minerals, pollutants, and --- critically --- radioactive isotopes.

In the years surrounding Project Argus, atmospheric nuclear testing deposited cesium-137 and strontium-90 across the Northern Hemisphere. The isotopes fell with rain and snow. They settled on soil, on water, on leaves, and on \textit{Cladonia stellaris}. The lichen absorbed them because the lichen absorbs everything. Caribou ate the lichen because caribou have eaten \textit{Cladonia} for ten thousand years --- it comprises up to seventy percent of their winter diet. The isotopes moved from lichen to caribou to the Inupiat and Gwich'in people who depend on caribou. A single nuclear test in Nevada or Kazakhstan or the South Atlantic became, through the patient chemistry of a star-tipped lichen, a dose of cesium in a woman's milk on the North Slope of Alaska.

The lichen did not know what it was accumulating. It had no discriminator, no threshold, no counter. It was not Explorer 4 with its scintillation crystals and its calibrated energy bins. It was a passive collector --- an organism that grew at two millimeters per year and recorded everything the atmosphere delivered, in the same way that a tree ring records drought and fire and plenty. What the Argus tests created in the sky, the lichen received on the ground. What Explorer 4 measured with instruments, \textit{Cladonia} measured with its body.

\begin{center}
$\bullet\quad\bullet\quad\bullet$
\end{center}

Carl Jung died on June 6, 1961, three years after the Argus tests. He had spent his life studying the parts of the self that remain invisible to the self --- the shadow, the anima, the collective unconscious. He believed that what we cannot see still shapes us. That the unconscious is not absence but presence: a living structure, as real as bone, that determines behavior from beneath the threshold of awareness.

The radiation belts are Jung's shadow made physical. They surround the Earth in an invisible shell. You cannot see them, feel them, taste them, or hear them. They interact with nothing that human senses detect. And yet they are there --- a structure of trapped energy, shaped by the magnetic field into a geometry as precise and as invisible as the unconscious. The Van Allen belts existed for four and a half billion years before Explorer 1 discovered them in January 1958. They shaped the space environment the entire time. Spacecraft that ventured into them would have had their electronics damaged, their crews irradiated, their missions compromised --- and nobody would have known why, because the belts were invisible.

Explorer 4 made the shadow visible. Not the natural shadow --- Van Allen had already done that --- but the artificial one. The shadow we created. Project Argus proved that humans could modify the radiation environment of the entire planet by detonating three small nuclear weapons at the right altitude. The belts we made were temporary. The belts that were already there are permanent. Both are invisible. Both are real. Both shape the environment for everything that passes through them.

The lichen knows. Not in the way that Jung knew, with language and symbol and analytic framework. In the way that a body knows: by absorbing, by accumulating, by carrying the evidence in its tissue where anyone with the right instrument can read it. Cut a cross-section of \textit{Cladonia stellaris} from the 1958 growth layer and run a gamma spectrometer across it. The cesium peak is there. The shadow left a mark.

\begin{center}
$\bullet\quad\bullet\quad\bullet$
\end{center}

Kohler finished his tour on the Norton Sound and went home to Modesto. He married a woman named Linda who taught third grade. He worked at the Gallo winery for thirty-two years. He never told anyone about the green light in the South Atlantic because the project was classified until 1959, and by then nobody asked. When his grandchildren found his Navy photographs in a shoebox in the garage, they asked him what the ship with the rocket was. He said it was a weapons test. They asked if it was scary. He said no. He said it was beautiful.

He said the sky turned green, and the green moved like a curtain, and for five minutes in the middle of the South Atlantic on a night in August 1958, the invisible became visible, and what he saw was not a weapon but a light, and the light was the color of something that had no name, and it faded the way all light fades, slowly and then all at once, until the sky was black again and the stars came back and the water was flat and calm and without feature, the way it had been before they put a hole in the magnetosphere and filled it with fire.

The lichen is still growing. Two millimeters a year. Absorbing whatever falls.

\end{document}
